% This file was converted to LaTeX by Writer2LaTeX ver. 1.9.9
% see http://writer2latex.sourceforge.net for more info
\documentclass{article}
\usepackage{calc,amsmath,amssymb,amsfonts}
\usepackage[T1]{fontenc}
\usepackage[italian]{babel}
\usepackage[style=numeric,backend=biber]{biblatex}
\author{HP}
\date{2026-07-26}
\begin{document}
\section{Prefazione}
Ci sono progetti che nascono, crescono e si concludono nel giro di pochi mesi.

Altri, invece, rimangono silenziosamente con noi per gran parte della vita.

Questo libro appartiene alla seconda categoria.

L'idea di realizzare un videogioco per Commodore 64 non è nata mentre scrivevo queste pagine. Non è nata neppure quando
ho iniziato a lavorare come programmatore. È molto più antica. Appartiene agli anni in cui un computer a otto bit
sembrava capace di aprire mondi sconosciuti e ogni nuova istruzione del BASIC dava l'impressione di avvicinarsi un po'
di più a qualcosa che sembrava quasi irraggiungibile.

Come molti ragazzi della mia generazione, anch'io sognavo di creare un videogioco.

Non una semplice prova di programmazione, ma un vero videogioco. Uno di quelli che, allora, osservavamo con meraviglia
cercando di immaginare come fosse possibile far muovere uno sprite, rilevare una collisione, far comparire un effetto
sonoro o gestire decine di oggetti con una macchina che oggi definiremmo estremamente limitata.

Quel sogno, però, è rimasto incompiuto.

La vita ha preso altre strade. Lo studio, il lavoro, i progetti professionali e le responsabilità hanno progressivamente
sostituito quelle lunghe giornate trascorse davanti al Commodore 64. Eppure, senza che me ne rendessi conto, qualcosa
di quel periodo è rimasto sempre presente.

Non il desiderio di tornare indietro.

Non la nostalgia.

Piuttosto la curiosità di capire, con gli strumenti e l'esperienza maturati negli anni, come sarebbe stato affrontare
oggi quella stessa sfida.

È proprio da questa domanda che nasce questo libro.

Non dall'ambizione di scrivere un manuale sul BASIC del Commodore 64 e nemmeno dal desiderio di raccontare alcuni
ricordi personali. Entrambe queste componenti sono presenti, ma nessuna delle due rappresenta il vero motivo per cui
queste pagine esistono.

Il vero obiettivo è un altro.

Vorrei accompagnare il lettore nella realizzazione di un videogioco così come avrebbe potuto nascere negli anni Ottanta:
un problema alla volta, una scelta alla volta, un errore alla volta. Senza scorciatoie. Senza soluzioni calate
dall'alto. Cercando di ricostruire non soltanto il codice, ma soprattutto il modo di ragionare che rendeva possibile
costruire software su una macchina con risorse tanto limitate.

Per questo motivo il libro non segue la struttura di un manuale tradizionale.

Gli argomenti non compaiono perché è arrivato il momento di spiegarli, ma perché il videogioco ne rende necessaria
l'introduzione. Ogni nuova conoscenza nasce da un problema reale, ogni soluzione modifica il progetto e ogni capitolo
rappresenta un passo avanti nella costruzione del programma.

Accanto a questo percorso tecnico ne scorre un altro, più discreto.

È il racconto delle persone, degli incontri e delle esperienze che hanno acceso quella passione e che, molti anni dopo,
hanno reso possibile trasformarla in un'opera compiuta.

Le due dimensioni procedono insieme perché, almeno per me, non è possibile separarle. Dietro ogni programma esistono
sempre persone, maestri, errori, intuizioni e occasioni che finiscono per influenzare il modo in cui impariamo a
pensare.

Questo libro è nato anche grazie a un lungo lavoro di progettazione e revisione, nel quale mi sono avvalso, tra gli
altri strumenti, anche del confronto con un'intelligenza artificiale. Il suo contributo è stato quello di un
interlocutore con cui discutere idee, verificare la coerenza dell'opera e affinare molte scelte editoriali. Le
riflessioni, i ricordi, le decisioni e ogni responsabilità riguardo ai contenuti appartengono però esclusivamente a me.

Non so quale sarà il risultato di questo percorso per ciascun lettore.

Mi auguro soltanto che, arrivato all'ultima pagina, non abbia semplicemente imparato qualcosa in più sul Commodore 64.

Spero che abbia avuto la sensazione di aver partecipato davvero alla nascita di un videogioco.

Perché, in fondo, è questa l'esperienza che avrei voluto vivere quando tutto ebbe inizio.

E che, molti anni dopo, ho finalmente avuto l'occasione di raccontare.


\bigskip
\end{document}
